\begin{lemma}[Smooth compactly supported functions are dense in $L^1(\gamma_n)$]
Fix $n \in \mathbb{N}_0$, let $\lambda$ be Lebesgue measure on $\mathbb{R}^n$ and
let $\gamma_n$ be the Gaussian measure of \noderef{GaussMeasure}.  Let
$g \in L^1(\gamma_n)$, and write $g$ also for its canonical representative, a
function $\mathbb{R}^n \to \mathbb{R}$.  Let $\delta > 0$.  Then there is a
function $\varphi : \mathbb{R}^n \to \mathbb{R}$ which is smooth, has compact
support and is $\gamma_n$-integrable, such that
\[
  \|g - \varphi\|_{L^1(\gamma_n)}
    = \int_{\mathbb{R}^n}|g(y) - \varphi(y)|\,d\gamma_n(y) < \delta ,
\]
the integral being read in $[0,+\infty]$ as the $L^1(\gamma_n)$ seminorm of the
difference, so that the inequality asserts in particular that it is finite.

This is the density statement the reference paper uses when it writes ``choose
$f_k \in C_c^\infty$ with $f_k \to f$ in $L^1(\gamma)$'': applying the lemma
with $\delta = 1/(k+1)$ produces exactly such a sequence.  It is stated plainly,
with no side condition on the support of $\varphi$ beyond compactness, because
that is the form the paper uses and the form that is reusable; callers that need
$\varphi$ to vanish on a prescribed ball obtain that for themselves by
multiplying by a cut-off, which costs nothing because the function being
approximated already vanishes there.

\emph{Why the tablet's existing approximation results do not give this.}  Two
of them are nearby and neither suffices, which is why this statement is
separate from both; they are named here for contrast only, are not used in the
proof below, and are correspondingly neither cited as dependencies nor
imported.  \texttt{L1ApproxByL2Truncations} produces elements of
$L^2(\gamma_n)$ whose $L^1(\gamma_n)$ images converge to $g$; those are
equivalence classes carrying no smoothness and no support control whatever.
\texttt{MollifierL2Local} concludes convergence of mollifications in $L^2$ of
\emph{Lebesgue} measure restricted to a compact set --- a local statement, about
a different measure and a different exponent.  Neither yields a $C_c^\infty$
approximant in $L^1(\gamma_n)$.
\end{lemma}
\begin{proof}
Write $\varrho(y) = \pi^{-n/2}e^{-|y|^2}$, which is smooth, strictly positive and
bounded above by $\pi^{-n/2}$ at every point of $\mathbb{R}^n$.  By
\noderef{GaussMeasure}, $\gamma_n$ is $\lambda$ weighted by $\varrho$, so that
for every measurable $w : \mathbb{R}^n \to \mathbb{R}$
\begin{equation*}
  \int_{\mathbb{R}^n}|w|\,d\gamma_n
    = \int_{\mathbb{R}^n}|w(y)|\,\varrho(y)\,d\lambda(y) , \tag{1}
\end{equation*}
both sides read in $[0,+\infty]$.

\emph{Step 0: measurability bookkeeping.}
The canonical representative $g$ is $\gamma_n$-almost everywhere equal to a Borel
function $\widetilde g$.  By clause (4) of \noderef{GaussLebesgueBallComparison}
the measures $\gamma_n$ and $\lambda$ have the same null sets, so
$g = \widetilde g$ also $\lambda$-almost everywhere.  Replacing $g$ by
$\widetilde g$ changes neither side of the asserted inequality --- the
integrand $|g - \varphi|$ changes only on a $\gamma_n$-null set --- so we may and
do assume $g$ Borel measurable.  Being the representative of an element of
$L^1(\gamma_n)$, it satisfies $\int|g|\,d\gamma_n < \infty$.

\emph{Step 1: transport to $L^1(\lambda)$.}
Put $h(y) = g(y)\varrho(y)$, a Borel function.  By $(1)$,
$\int|h|\,d\lambda = \int|g|\,d\gamma_n < \infty$, so $h \in L^1(\lambda)$.

\emph{Step 2: a continuous compactly supported approximant to $h$.}
By the density of compactly supported continuous functions in $L^p$ of a regular
Borel measure (shared context, applied to $\mu = \lambda$, $p = 1$ and the
function $h$) there is $u : \mathbb{R}^n \to \mathbb{R}$ which is continuous,
has compact support, lies in $L^1(\lambda)$ and satisfies
\begin{equation*}
  \int_{\mathbb{R}^n}|h - u|\,d\lambda \leq \frac{\delta}{3} . \tag{2}
\end{equation*}
The same shared-context item records that such a $u$ is bounded and uniformly
continuous on $\mathbb{R}^n$.

\emph{Step 3: a smooth compactly supported approximant to $u$.}
For $\varepsilon > 0$ fix $\rho_\varepsilon : \mathbb{R}^n \to \mathbb{R}$ which
is smooth, nonnegative, has $\operatorname{supp}\rho_\varepsilon \subseteq
\overline B(0,\varepsilon)$, and satisfies $\int\rho_\varepsilon\,d\lambda = 1$.
Such a family exists.  For $n \geq 1$: take the standard radial bump $b$ at
centre $0$ with inner radius $\varepsilon/2$ and outer radius $\varepsilon$, so
that $b$ is smooth, takes values in $[0,1]$, equals $1$ on
$\overline B(0,\varepsilon/2)$ and is supported in $B(0,\varepsilon)$; then
$\int b\,d\lambda > 0$, because $b$ is nonnegative, continuous and equal to $1$
on a ball of positive Lebesgue measure, and $\rho_\varepsilon = b/\int b\,d\lambda$
has all four properties.  For $n = 0$ the space $\mathbb{R}^0$ is a single point
of unit $\lambda$-mass and $\rho_\varepsilon \equiv 1$ has all four properties.  Set
\[
  u_\varepsilon(x) = (\rho_\varepsilon * u)(x)
    = \int_{\mathbb{R}^n}\rho_\varepsilon(t)\,u(x-t)\,d\lambda(t) .
\]
By the shared-context item on smoothness of a mollification, $u_\varepsilon$ is
defined at every $x$ and is smooth.

\emph{Support.}  Let $S = \operatorname{supp}u$, compact, and let
$K = \{x : \operatorname{dist}(x,S) \leq 1\}$, which is closed and bounded,
hence compact, and has $\lambda(K) < \infty$ (shared context).  If
$0 < \varepsilon \leq 1$ and $x \notin K$ then for every $t$ with
$|t| \leq \varepsilon$ we have $x - t \notin S$, so $u(x-t) = 0$; since
$\rho_\varepsilon$ vanishes off $\overline B(0,\varepsilon)$, the integrand
above vanishes identically and $u_\varepsilon(x) = 0$.  Thus
$\{u_\varepsilon \neq 0\} \subseteq K$, and as $K$ is closed,
$\operatorname{supp}u_\varepsilon \subseteq K$ is compact.

\emph{Uniform convergence.}  Since $\int\rho_\varepsilon\,d\lambda = 1$ we have
$u(x) = \int\rho_\varepsilon(t)u(x)\,d\lambda(t)$, so for every $x$
\[
  |u_\varepsilon(x) - u(x)|
    = \Bigl|\int\rho_\varepsilon(t)\bigl(u(x-t)-u(x)\bigr)d\lambda(t)\Bigr|
    \leq \sup_{|t|\leq\varepsilon}\ \sup_{w\in\mathbb{R}^n}|u(w-t)-u(w)| ,
\]
using $\rho_\varepsilon \geq 0$, $\int\rho_\varepsilon = 1$ and that
$\rho_\varepsilon$ vanishes for $|t| > \varepsilon$.  By the uniform continuity
of $u$ (Step 2), given $\eta > 0$ there is $\varepsilon \in (0,1]$ with
$|u(w-t)-u(w)| \leq \eta$ for all $w$ and all $|t| \leq \varepsilon$, and then
$\sup_x|u_\varepsilon(x)-u(x)| \leq \eta$.  Since $u_\varepsilon - u$ vanishes
off $K$,
\[
  \int_{\mathbb{R}^n}|u_\varepsilon - u|\,d\lambda \leq \eta\,\lambda(K) .
\]
Choose $\eta > 0$ with $\eta\,\lambda(K) \leq \delta/3$ --- possible because
$\lambda(K) < \infty$ --- and let $v = u_\varepsilon$ for the corresponding
$\varepsilon$.  Then $v$ is smooth with compact support and
\begin{equation*}
  \int_{\mathbb{R}^n}|u - v|\,d\lambda \leq \frac{\delta}{3} . \tag{3}
\end{equation*}

\emph{Step 4: transport back, and conclusion.}
Define
\[
  \varphi(y) = \pi^{n/2}e^{|y|^2}\,v(y) ,
  \qquad\text{so that}\qquad
  \varphi(y)\,\varrho(y) = v(y) \text{ for every } y .
\]
The function $y \mapsto \pi^{n/2}e^{|y|^2}$ is smooth and nowhere zero, so
$\varphi$ is smooth and $\{\varphi \neq 0\} = \{v \neq 0\}$, whence
$\operatorname{supp}\varphi = \operatorname{supp}v$ is compact.  Being
continuous with compact support, $\varphi$ is bounded; and $\gamma_n$ is a
probability measure, hence finite (\noderef{GaussIsProb}), so a bounded
measurable function is $\gamma_n$-integrable: $\varphi \in L^1(\gamma_n)$.

Finally, by $(1)$ applied to $w = g - \varphi$, then
$(g(y)-\varphi(y))\varrho(y) = h(y) - v(y)$, then the triangle inequality in
$L^1(\lambda)$ together with $(2)$ and $(3)$,
\[
  \int_{\mathbb{R}^n}|g - \varphi|\,d\gamma_n
    = \int_{\mathbb{R}^n}|h - v|\,d\lambda
    \leq \int_{\mathbb{R}^n}|h - u|\,d\lambda
      + \int_{\mathbb{R}^n}|u - v|\,d\lambda
    \leq \frac{\delta}{3} + \frac{\delta}{3} < \delta .
\]
This is the assertion.
\end{proof}
